\section{重庆多联机系统设计}

本章按三层办公楼房间级参数完成重庆工况下的夏季冷负荷计算和VRF设备选型。重庆属于夏热冬冷高湿地区，设计重点为高温高湿叠加，峰值冷量和除湿压力最大。计算模型采用50个房间级热源参数和15个天正暖通计算分区，房间热源参数包括人员、照明、设备、新风和外窗朝向。

\subsection{室外设计参数与围护结构}

重庆室内设计参数统一取\SI{26}{\degreeCelsius}、相对湿度60\%，基准新风量40~m$^3$/(h$\cdot$人)。围护结构采用本设计的优化参数，和参考设计相比，南方城市强化遮阳控制，北方城市强化保温控制。

\begin{table}[!htbp]
\centering
\caption{重庆室外设计参数与围护结构参数}
\label{tab:chongqing-envelope}
\zihao{-5}
\begin{tabular}{@{}lccccccc@{}}
\toprule
城市 & 干球温度 & 湿球温度 & 日平均温度 & 外墙K & 屋面K & 外窗K & SC \\
 & \si{\degreeCelsius} & \si{\degreeCelsius} & \si{\degreeCelsius} & \si{W/(m^2\cdot K)} & \si{W/(m^2\cdot K)} & \si{W/(m^2\cdot K)} & -- \\
\midrule
重庆 & 35.5 & 27.6 & 30.7 & 0.72 & 0.52 & 2.50 & 0.34 \\
\bottomrule
\end{tabular}
\end{table}

\subsection{房间级冷负荷计算结果}

重庆全楼峰值冷负荷为\SI{150.48}{kW}，峰值冷负荷指标为\SI{110.8}{W/m^2}，峰值时刻为14:00。其中室内侧负荷为\SI{79.93}{kW}，新风冷负荷为\SI{70.55}{kW}，新风占比约46.9\%。主要高负荷房间见表~\ref{tab:chongqing-top-rooms}。

\begin{table}[!htbp]
\centering
\caption{重庆主要房间峰值冷负荷与热源参数}
\label{tab:chongqing-top-rooms}
\zihao{-5}
\begin{tabular}{@{}cccccccc@{}}
\toprule
楼层 & 房间编号 & 房间名称 & 面积 & 人数 & 新风量 & 峰值冷负荷 & 新风冷负荷 \\
 & -- & -- & \si{m^2} & 人 & \si{m^3/h} & \si{kW} & \si{kW} \\
\midrule
1 & 1001 & 外贸办公室 & 78.5 & 14 & 630 & 11.99 & 6.13 \\
2 & 2008 & 办公大厅 & 54.9 & 10 & 450 & 8.78 & 4.38 \\
3 & 3007 & 会客室 & 37.8 & 11 & 440 & 8.04 & 4.28 \\
3 & 3013 & 小型会议室 & 38.6 & 11 & 440 & 7.50 & 4.28 \\
1 & 1012 & 值班室 & 54.1 & 8 & 320 & 7.06 & 3.11 \\
1 & 1008 & 陈列室 & 73.9 & 7 & 245 & 6.76 & 2.38 \\
1 & 1011 & 接待大厅 & 61.0 & 7 & 245 & 6.17 & 2.38 \\
2 & 2001 & 办公室1 & 39.3 & 6 & 240 & 5.60 & 2.33 \\
3 & 3018 & 办公茶歇区 & 25.9 & 7 & 280 & 4.71 & 2.73 \\
3 & 3017 & 办公室4 & 36.1 & 5 & 200 & 4.54 & 1.95 \\
\bottomrule
\end{tabular}
\end{table}

\subsection{VRF室内外机选型与容量配比}

室内机按房间峰值冷负荷乘1.3安全系数选型，普通办公、会议和会客房间优先采用四面出风嵌入式室内机，卫生间、走廊、档案和储藏类房间采用薄型风管或壁挂式末端。重庆共配置48台室内机，室内机名义制冷量合计\SI{243.1}{kW}。每层设置一个室外机系统，容量配比范围为114.0\%--117.9\%，满足JGJ~174-2010中50\%--130\%的组合率要求。

\begin{table}[!htbp]
\centering
\caption{重庆各楼层VRF室内外机容量配比}
\label{tab:chongqing-vrf-ratio}
\zihao{-5}
\begin{tabular}{@{}ccccccc@{}}
\toprule
楼层 & 室内机台数 & 室内容量 & 室外机型号 & 室外机台数 & 室外容量 & 配比 \\
 & 台 & \si{kW} & -- & 台 & \si{kW} & \% \\
\midrule
1 & 12 & 70.1 & VRF-O615 & 1 & 61.5 & 114.0 \\
2 & 16 & 78.7 & VRF-O680 & 1 & 68.0 & 115.7 \\
3 & 20 & 94.3 & VRF-O800 & 1 & 80.0 & 117.9 \\
\bottomrule
\end{tabular}
\end{table}

从负荷组成看，重庆的新风冷负荷占比为46.9\%，该比例直接影响DOAS或新风预处理设备的容量。若施工图阶段采用厂家样本深化，应优先复核高负荷房间的室内机布置、冷媒管长修正和室外机通风散热条件。
